% English translation of checked_pages/esq_119_047.tex; source: cote 119, site 47, PDF page 48, typescript page 45.
% AI-translated by GPT-6 Astra (OpenAI, ChatGPT Pro mode).
\scanpage{119/47}
\mspage{45}

This work is not entirely isolated in relation to
my more recent interests. For example, my reflection on the
moduli multiplicities $M^{\wedge}_{g,\nu}$ and their stratified structure
revived reflection on a van Kampen theorem of dimension
$>$ 1 (also one of the favorite themes of the Bangor
group), and may have helped prepare the ground for
the work of greater scope the following year. The latter
also connects at times with a reflection dating from the same
year 1975 (or the following year) on a ``De Rham complex with
divided powers'', which was the subject of my last public
lecture, at IHES\tnote{TN-047-01}{IHES is the Institut des Hautes Études Scientifiques; see the \href{https://www.ihes.fr/en/}{institute's website}.} in 1976, and whose manuscript, entrusted
to someone I no longer remember after the lecture, is moreover lost.
It was at the time of this reflection that the intuition
of a ``schematization'' of homotopy types also began to germinate, which,
seven years later, I am trying to make precise in a (particularly
hypothetical) chapter of History of Models.

The work of reflection undertaken in Pursuing
Stacks is a little like a debt I would be discharging
toward a scientific past in which, for some fifteen
years (between 1955 and 1970), the development of cohomological
tools was the constant leitmotif in my work on the
foundations of algebraic geometry. If the present resumption of this
theme has taken on unexpected dimensions, it is not,
however, out of pity for a past, but because of the many unforeseen
things constantly breaking in, mercilessly upsetting the
plans and intentions envisaged --- a little like a tale from the Thousand and
One Nights, in which one's attention is kept in suspense through
twenty other tales before learning the final word of the first.


8. I have as yet said very little about the more down-to-earth
reflections on two-dimensional topological geometry, associated
in particular with my activities as a teacher and those called ``research
supervision''. On several occasions, I have seen\unsure{.}\tnote{TN-047-02}{The doubtful dot after ``vu'' is retained after the corresponding ``seen''.} opening up before

\margin{47}
